\chapter{Theoretical Compiler}
\label{chap:compiler}

In this chapter, we describe a theoretical compiler from \langc\ to \lange. The
idea to the compilation is simple, with special handling for dependent contract.
A monitor of flat contract will be converted into performing an effect. A
monitor of function contract will be recursively expanded similarly to the
reduction until only flat contracts exist. Monitor of dependent contract is
compiled similarly to the flat contract, with a slight difference for dependent
argument.

The compilation described above only add in effect performing, without their
handler component. Without a suitable handler and using handle expression, the
program cannot run. First, we define that each contract can be converted into a
handler. The handler accepts a non-function value, and check it against the
contract predicate. If the value passes the predicate, the program continues
(using the handler continuation) with the value. Else the program steps to a
error, which stops the program with the party responsible for the blame.

With a handler defined, we need to use it with a handle expression to ensure
effects performed are handled. The easiest approach is to apply the handle right
where the effect is performed with its corresponding handler. This does not give
any advantage to basic if-else compilation. Handlers can stay inside the same
context if they are used together. We suggest grouping them together, this
allows contracts to access the same context information.

Grouping them together is easy, and we can place a handle at the whole program.
However, this approach fails for dependent contract checking. Dependent contract
$\kappa_d$ checking happens while executing another contract $\kappa_c$. If
$\kappa_c$ is running inside a handler, which is true in our described
compilation, then the handler is moved into the continuation, leaving $\kappa_c$
running without a handle context. To resolve this, we have to handle $\kappa_d$
using a handler of all possible effects that might arise during execution of
$\kappa_d$. This loses the general context shared across all other contracts,
but guarantees all effects are handled.

The next section illustrates many examples aim to give an intuition to the
reader before the later formal section, where a compilation model between
\langc\ to \lange\ is presented.

\section{Examples}

We first examine the flat contract case in Figure
\ref{fig:example-compile-flat}. A flat contract compiles into performing effect
operation $\attrE{check}$, a unique effect name generated for this contract.
The handler $\attrE{h}$ defined with only the operation $\attrE{check}$, and
not applied immediately to the compiled expression. Handler $\attrE{h}$ will be
kept going forward the compilation process.

\begin{figure}
  \begin{align*}
    \cMon{k}{l}{j}{\cFlat{\cLam{x : \cTnum}{\attrC{zero?}(x \attrC{-}
  10)}}}{\attrC{9}}
  \end{align*}

$\Downarrow$

\begin{align*}
  \attrE{h} \ &\triangleq \ \attrE{\{} \attrE{check \to} \\
  &\qquad \eTAbs{\alpha}{
    \eLam{x : \alpha}{
      \eLam{kont : \alpha \attrE{\to \langle\rangle} \alpha}{
        \eIf{(\eLam{x:\attrE{num}}{\attrE{zero?(x - 10)}}) \ \attrE{x}}
            {\attrE{kont \ x}}
            {\eBlame{k}}
      }
    }
  }
  \attrE{\}} \\[2ex]
    &\ePerform{check}{\attrE{num}} \ \attrE{9}
\end{align*}

  \caption[Example compilation of flat contract]
  {Example compilation of flat contract\\
  \small
  }
\label{fig:example-compile-flat}
\end{figure}


A function contract is then straightforward by recursively compile the argument
and the application result into \lange. Figure
\ref{fig:example-compile-function} demonstrates the unrolling of function
contract into single flat contracts and compiling them into their corresponding
operation invocation. Similarly to flat contract, the group of handlers are not
used immediately.

A dependent contract compiles similarly to a function contract. However,
because the handler scope during checking the post contract is pushed into the
continuation, the dependent contract must be wrapped in its own handler scope.
Figure \ref{fig:example-compile-dependent} describe the procedure for a
\textit{basic type} argument. Because the argument is a basic type, there is
only one operation in the handler $h_3$. If the argument is a function, then
$h_3$ should capture all possible effects, collectable after compiling
recursively with known contract.


\begin{figure}
\begin{align*}
  \cMon{k}{l}{j}{\cKFunc{(\cKFunc{\cFlat{k_1}}{\cFlat{k_2}})}{\cFlat{k_3}}}{f}
\end{align*}
$\Downarrow$
\begin{align*}
  \cLam{g}{\cMon{k}{l}{j}{\cFlat{k_3}}{f \ (\cMon{l}{k}{j}{\cKFunc{\cFlat{k_1}}{\cFlat{k_2}}}{g})}}
\end{align*}
$\Downarrow$
\begin{align*}
  \cLam{g}{\cMon{k}{l}{j}{\cFlat{k_3}}{f \ (\cLam{x}{\cMon{l}{k}{j}{\cFlat{k_2}}{g \ (\cMon{k}{l}{j}{\cFlat{k_1}}{x})}})}}
\end{align*}
$\Downarrow$

\begin{align*}
  % 1. Define the handlers with their correct blame polarities
  \attrE{h_1} \ &\triangleq \ \attrE{\{} \attrE{check_{k_1} \to} \dots \eBlame{k} \attrE{\}} \\
  \attrE{h_2} \ &\triangleq \ \attrE{\{} \attrE{check_{k_2} \to} \dots \eBlame{l} \attrE{\}} \\
  \attrE{h_3} \ &\triangleq \ \attrE{\{} \attrE{check_{k_3} \to} \dots \eBlame{k} \attrE{\}} \\[2ex]
  % 2. Build the nested evaluations bottom-up
  \attrE{e_1} \ &\triangleq \ {\ePerform{check_{k_1}}{\tau_1} \ \attrE{x}} \\
  \attrE{e_2} \ &\triangleq \ {\ePerform{check_{k_2}}{\tau_2} \ (\attrE{g} \ \attrE{e_1})} \\
  \attrE{e_3} \ &\triangleq \ {\ePerform{check_{k_3}}{\tau_3} \ (\attrE{f} \ (\eLam{x}{\attrE{e_2}}))} \\[2ex]
  % 3. The final compiled expression
  &\eLam{g}{\attrE{e_3}}
\end{align*}

\caption{Example compilation of higher-order contract}
\label{fig:example-compile-function}
\end{figure}

\begin{figure}
\begin{align*}
  \cMon{k}{l}{j}{\cKDep{\cFlat{k_1}}{x}{\cFlat{k_2}}}{g}
\end{align*}

$\Downarrow$

\begin{align*}
  \cLam{x}{\cMon{k}{l}{j}{\cSubst{\cMon{l}{j}{j}{\cFlat{k_1}}{x}}{x}{\cFlat{k_2}}}{g \ (\cMon{l}{k}{j}{\cFlat{k_1}}{x})}}
\end{align*}

$\Downarrow$

\begin{align*}
  % 1. Define the handlers with their correct blame polarities
  \attrE{h_1} \ &\triangleq \ \attrE{\{} \attrE{check_{k_1} \to} \dots \eBlame{l} \attrE{\}} \\
  \attrE{h_2} \ &\triangleq \ \attrE{\{} \attrE{check_{k_2} \to} \\
  &\qquad \eTAbs{\alpha}{
    \eLam{y : \alpha}{
      \eLam{kont : \alpha \attrE{\to \langle\rangle} \alpha}{
        \eIf{k_2[x:=x_{dep}]}
            {\attrE{kont \ y}}
            {\eBlame{k}}
      }
    }
  } \\
  \attrE{h_3} \ &\triangleq \ \attrE{\{} \attrE{check'_{k_1} \to} \dots \eBlame{l} \attrE{\}} \\[2ex]
  % 2. Build the nested evaluations bottom-up
  \attrE{x_{arg}} \ &\triangleq \ {\ePerform{check_{k_1}}{\tau_1} \ \attrE{x}}\\
  \attrE{x_{dep}} \ &\triangleq \ {\eHandle{h_3}{\ePerform{check'_{k_1}}{\tau_1} \ \attrE{x}}}\\
  \attrE{e_1}     \ &\triangleq \ {\ePerform{check_{k_2}}{\tau_2} \ {(\attrE{g} \ \attrE{x_{arg}})}} \\[2ex]
  % 3. The final compiled expression
  &\eLam{x}{\attrE{e_1}}
\end{align*}

\caption{Example compilation of dependent contract}
\label{fig:example-compile-dependent}
\end{figure}


\section{Type-Directed Compilation}

The compilation relies heavily on the type of argument. Therefore, we provide a
type-directed compilation from \langc\ to \lange. Compilation relation
$\cCompile{\cGamma}{e}{\tau}{h}{e}$ takes an expression $\attrC{e}$ of type
$\attrC{τ}$ with context $\cGamma$ and compiles into $\attrE{e}$ with a handler
$\attrE{h}$.


The compilation is recursive, types are used to form well-typed \lange\ expressions. We define the following theorems for our compiler: Type Preservation, Observable Equivalance, and Blame Preservation. Type Preservation Theorem states that a well-typed program in \langc\ compiles into a well-typed program in \lange\ typed using the same basic type and has no unhandled effects. Observable Equivalance Theorem states that if the \langc\ program reduces to a value, then the compiled \lange\ program reduces to a value also. Blame Preservation Theorem states that if \langc\ program blames a party, then the compiled \lange\ program also blame the same party.

Because \lange\ has different types from \langc, we define $\symtt{T}$ as the type conversion from \cGamma\ to \eGamma. The conversion is simple because \langc\ only has basic type, so we map to the basic type of \lange.

In Figure \ref{fig:compilation}, we provide some interesting compilation rules, other expressions left out should be trivial. Types used in \lange\ that copied from \langc context are assumed to be converted.

\begin{figure}[htbp]
  \centering
  \textbf{Environment Translation}
  \begin{align*}
    \symtt{T}(\attrC{\cdot}) \ &\triangleq \ \attrE{\cdot} \\
    \symtt{T}(\cGamma, \attrC{x : \tau}) \ &\triangleq \ \symtt{T}(\cGamma), \attrE{x : \tau^*}
  \end{align*}

  \vspace{1em}
  \textbf{Type Translation ($\tau \mapsto \tau^*$)}
  \begin{align*}
    \attrC{b}^* \ &\triangleq \ \attrE{b} \qquad \text{(where } \attrC{b} \text{ is a basic type like } \attrC{num} \text{)} \\
    (\attrC{\tau_1 \to \tau_2})^* \ &\triangleq \ \attrE{\tau_1^* \to \langle\rangle\tau_2^*}
  \end{align*}

  \caption[Type and Environment Translation from \langc\ to \lange]
  {Type and Environment Translation from \langc\ to \lange\\ \small Basic types
  are converted into basic types in kind system. Function types are
  computation without any effects. Contract type has no equivalent type in
  \lange. Context type translation $\symtt{T}$ converts all variables in \langc\
  into their equivalent types in \lange.}
  \label{fig:env-translation}
\end{figure}

\begin{theorem}[Type Preservation]
  If $\cCompile{\cGamma}{e}{\tau}{\Sigma}{e'}$ then $\eTypesC{\eGamma}{e'}{τ^*}{\langle\rangle}$
  where $\symtt{T}(\cGamma) \subseteq \eGamma$
\end{theorem}

\begin{theorem}[Observable Equivalance]
  If $\attrC{e} \cStepMany \attrC{v}$ and $\cCompile{\cdot}{e}{\tau}{\Sigma}{e'}$ then there
  exists a $\attrE{v'}$ such that $\attrE{e'}\eStepMany \attrE{v'}$
\end{theorem}

\begin{theorem}[Blame Preservation]
  If $\attrC{e} \cStepMany \cErr{l}{p}$ and $\cCompile{\cdot}{e}{\tau}{\Sigma}{e'}$ then $\attrE{e'}\eStepMany \eBlame{p}$
\end{theorem}



\begin{figure}[htbp]
  \centering
  \begin{mathpar}
  \inferrule[Compile-Var]
    {\attrC{x} : \attrC{\tau} \in \cGamma}
    {\cCompile{\cGamma}{x}{\tau}{\{\}}{x}}

  \inferrule[Compile-Val]
    {\strut}
    {\cCompile{\cGamma}{v}{\tau}{\{\}}{v}}

  \inferrule[Compile-Lambda]
    {\cCompile{\cGamma, \attrC{x} : \attrC{\tau_1}}{e}{\tau_2}{h}{e'}}
    {\cCompile{\cGamma}{\cLam{x:\tau_1}{e}}{\tau_1 \to \tau_2}{h}{\eLam{x:\tau_1^*}{e'}}}

  \inferrule[Compile-If]
    { \cCompile{\cGamma}{e_1}{bool}{h_1}{e_1'} \\
      \cCompile{\cGamma}{e_2}{\tau}{h_2}{e_2'} \\
      \cCompile{\cGamma}{e_3}{\tau}{h_3}{e_3'} }
    { \cCompile{\cGamma}{\cIf{e_1}{e_2}{e_3}}{\tau}{h_1 \cup h_2 \cup h_3}{\eIf{e_1'}{e_2'}{e_3'}} }

  \inferrule[Compile-BinOp-Nat]
    { \cCompile{\cGamma}{e_1}{nat}{h_1}{e_1'} \\
      \cCompile{\cGamma}{e_2}{nat}{h_2}{e_2'} }
    { \cCompile{\cGamma}{e_1 \ \attrC{\oplus_{nat}} \ e_2}{nat}{h_1 \cup h_2}{e_1' \ \attrE{\oplus_{nat}} \ e_2'} }

  \inferrule[Compile-BinOp-Bool]
    { \cCompile{\cGamma}{e_1}{bool}{h_1}{e_1'} \\
      \cCompile{\cGamma}{e_2}{bool}{h_2}{e_2'} }
    { \cCompile{\cGamma}{e_1 \ \attrC{\oplus_{bool}} \ e_2}{bool}{h_1 \cup h_2}{e_1' \ \attrE{\oplus_{bool}} \ e_2'} }

  \inferrule[Compile-App]
    {\cCompile{\cGamma}{e_1}{\tau_1 \to \tau}{h_1}{e_1'} \\
     \cCompile{\cGamma}{e_2}{\tau_1}{h_2}{e_2'}}
    {\cCompile{\cGamma}{e_1 \ e_2}{\tau}{h_1 \cup h_2}{e_1' \ e_2'}}

  \inferrule[Compile-MonFlat]
    { \cCompile{\cGamma}{e}{\tau}{h}{e'} \\
      \attrE{check_{\kappa} ∉  h} \\
      \attrE{h_\kappa} \triangleq \attrE{\{} \attrE{check_{\kappa} \to f} \attrE{\}} \\
      \attrE{f} \triangleq
        \eTAbs{\alpha}{ \eLam{x : \alpha}{ \eLam{kont :
          \alpha \attrE{\to \langle\rangle} \alpha}{ \eIf{\kappa \ x} {\attrE{kont \ x}}
          {\eBlame{k}} } } } }
      {
        \cCompile{\cGamma}{\cMon{k}{l}{j}{\cFlat{\kappa}}{e}}{\tau}{h
        \cup h_\kappa}{\ePerform{check_{\kappa}}{\tau^*}\ e'} }

  \inferrule[Compile-MonFunc]
    { \attrC{x} ∉ \cGamma \\
      \cCompile{\cGamma, \attrC{x} : \attrC{\tau_1}}{\cMon{k}{l}{j}{\kappa_2}{e
      \ (\cMon{l}{k}{j}{\kappa_1}{x})}}{\tau_2}{h}{e'} }
    { \cCompile{\cGamma}{\cMon{k}{l}{j}{\kappa_1 \to \kappa_2}{e}}{\tau_1 \to
    \tau_2}{h}{\eLam{x:\tau_1^*}{e'}} }

  \inferrule[Compile-MonDep]
    { \attrC{y},\delta_y ∉ \cGamma \\
      \cCompile{\cGamma, \attrC{y} :
      \attrC{\tau_1}}{\cMon{l}{j}{j}{\kappa_1}{y}}{\tau_1}{h_1}{e_y} \\
      \cCompile{\cGamma, \attrC{y} :
      \attrC{\tau_1}}{\cMon{k}{l}{j}{\{{\color{black}\delta_y}/ x\}\kappa_2}{e \
      (\cMon{l}{k}{j}{\kappa_1}{y})}}{\tau_2}{h_2}{e'} \\
      \attrE{h'_2} = \attrE{h_2}[\delta_y\coloneqq{\symtt{wrap}(\attrE{h_1},\attrE{e_y})}]}
      {
        \cCompile{\cGamma}{\cMon{k}{l}{j}{\cKDep{\kappa_1}{x}{\kappa_2}}{e}}{\tau_1
        \to \tau_2}{h_1 \cup h'_2}{\eLam{y:\tau_1^*}{e'}} }
  \end{mathpar}

  \textbf{Helper function}
  \begin{align*}
    \symtt{wrap}(\attrE{h},\attrE{e}) \ &\triangleq \ \eHandle{h}{e}, \attrE{e} \neq \eLam{x}{e} \\
    \symtt{wrap}(\attrE{h},\eLam{x}{e}) \ &\triangleq \ \eLam{x}{\eHandle{h}{e}}
  \end{align*}



  \caption[Type-Directed Compilation from \langc\ to \lange]
  {Type-Directed Compilation from \langc\ to \lange\\
  \small
  Original types in \langc\ is converted into their equivalent types in \lange\ after compilation. When
  compiling dependent contracts, a separate variable is used to store the dependent
  argument then replaced on top the compiled expression.}
  \label{fig:compilation}
\end{figure}

